\chapter{Polyakov Action, T Duality, and Compactification of the HMT State in M-Theory Branes}
\label{chap:teoria-m}

\section{Extended persistence and world surface}

A physical string is an extended one-dimensional object whose evolution sweeps out
a worldsheet. Each point of the string has an internal coordinate,
and each instant adds a temporal coordinate; their combination forms
the domain on which the induced metric, tension, and background fields
live. A closed string has a circular spatial section. An open string
has endpoints subject to boundary conditions. Its vibrations organize
states, and its tension converts the area traversed by the worldsheet into
action.

A brane of dimension \(p\) is an extended object whose history forms a
worldvolume of dimension \(p+1\). Branes receive string endpoints,
carry fields on their volume, and may wrap cycles of a compact
space. This definition fixes three different data: the dimension of the object,
the dimension of its history, and the conditions governing its boundary. The
tension, charges, and spectrum belong to the dynamical realization that
completes those geometric data.

HMT reaches this boundary from an architecture of persistent routes. A
route preserves the orientation, sheet, phase, quotient, carry, and memory of the
crossings that have prolonged it. The passage to an extended object assigns that
genealogy to a section distributed over an internal coordinate. The
composition of routes is then effected by sewing segments; the boundary
memory records windings and intersections; the action measures the
complete persistence over the traversed surface. This map preserves the
causal order: first the enriched state is constructed, and its realization
as a string, brane, or field is declared thereafter.

The Polyakov action provides a precise form for the physical codomain:
\[
 S_{\rm P}[X,h]
 =-\frac{T}{2}\int d^2\sigma\,
 \sqrt{-h}\,h^{ab}g_{MN}(X)
 \partial_aX^M\partial_bX^N.
\]
Here \(T\) is the tension, \(h_{ab}\) the worldsheet metric and
\(X^M\) the map into the target space. An HMT string realization must
produce these objects, their state space and their symmetries from a typed
map; the route structure supplies the genealogical domain on which
that map acts.

\section{Exact screens in dimensions \(11\), \(10\), \(26\) and \(24\)}

The first screen starts from the extended ternary code
\(C_W=[12,6,6]_3\). Puncturing one coordinate produces the code
\(G_{11}=[11,6,5]_3\). Its Hamming balls of radius two have cardinality
\[
 V_3(11,2)=1+22+220=243=3^5,
\]
and the \(3^6\) codewords satisfy
\[
 3^6V_3(11,2)=3^{11}.
\]
Therefore, these balls partition the ambient space exactly. The integer
eleven here measures the length of a perfect code and retains its
combinatorial type.

A second screen uses the real permutation module \(V_{12}\) associated with
the twelve icosahedral vertices. The complement of the uniform vector
decomposes as
\[
 V_{11}=V_3\oplus V_{3'}\oplus V_5.
\]
The dodecaphase state \(K\) polarizes a line \(\ell_K\subset V_3\) and
determines the projector
\[
 P_{10}=P_{11}-\frac{u_Ku_K^{\mathsf T}}{\lVert u_K\rVert^2}.
\]
From it descend the exact decompositions.
\[
 V_{11}=\ell_K\oplus V_{10},\qquad
 V_{10}=V_{3'}\oplus W_7(K),
\]
with dimensions \(11=1+3+7\) and \(10=3+7\). The step
\(12\to11\) eliminates the uniform mode; the step \(11\to10\) eliminates the
direction marked by \(K\). A physical compactification adds a
temporal line, a stable rank-six lattice, and the transport that preserves the
three extended directions.

The third screen starts from the lattice
\[
 L_{26}=\Lambda_{24}\oplus U\cong II_{25,1},
\]
where \(U\) is the integral hyperbolic plane. For a primitive isotropic
vector \(e_+\), the reduction
\[
 e_+^\perp/\mathbb Ze_+\cong\Lambda_{24}
\]
is exact. The orthogonality condition removes one direction and the quotient
removes the null direction contained in its orthogonal complement; the result is rank
twenty-four. Code, real module, and lattice belong to distinct
categories. Their explicit maps make it possible to relate their ranks while preserving the
provenance of each.

The incidence between the six faces and the eight vertices of the cube also provides
a screen of spatial and temporal signature. If \(M\) is its incidence
matrix, then
\[
 MM^{\mathsf T}=12P_u+4P_-.
\]
The quotient by the kernel has rank four and decomposes as
\(1\oplus3\). Upon declaring the uniform line to be the temporal direction, the
normalized form becomes \(\operatorname{diag}(-1,1,1,1)\). The incidence proves
the rank and the spectral ratio; the declared orientation fixes the signature.

\section{Moment, winding and reciprocal radius}

In a compact coordinate of radius \(R\), the integer \(m\) records momentum and
the integer \(w\) records winding. The form
\[
 E_R(m,w)=\frac{m^2}{R^2}+w^2R^2
\]
has the involution
\[
 \mathsf T(m,w,R)=(w,m,R^{-1}),
 \qquad
 E_R(m,w)=E_{R^{-1}}(w,m).
\]
In \(\ell^2(\mathbb Z^2)\), the exchange
\(U_T\lvert m,w\rangle=\lvert w,m\rangle\) provides the unitary
equivalence
\[
 U_TH(R)U_T^{-1}=H(R^{-1}).
\]
This is the algebraic residence of radius duality: momentum and
winding are two reciprocal readings of the same lattice sector. The
realization as T-duality incorporates the string scale, the oscillators and
the action on the physical state space. At the self-dual radius both
readings share the same scale and the exchange acts within the same
spectrum.

\section{Interface in dimension \(11\): spectral triple, supersymmetry and operator \(\mathcal Q^2=H\)}

M-theory organizes an eleven-dimensional physical codomain with metric
\(g_{MN}\), three-form \(C_{MNP}\), gravitino \(\psi_M\), action, and
supersymmetry transformations. The bridge is written as a map
\[
 \mathfrak R_M:\mathcal X_{\rm HMT}^{\rm enr}
 \longrightarrow\mathcal X_{11}^{\rm phys}.
\]
Its source preserves phase, orientation, sewing and memory; its codomain adds
the fields and the eleven-dimensional dynamics. The compactification of one
direction must transport these fields to a ten-dimensional screen,
preserve the low-energy limit and intertwine the declared dualities.

A candidate compactification is typed by a triple
\((x_\infty,\kappa,\mathcal U)\). Section \(x_\infty\) represents a
prolongable persistence; \(\kappa\) carries it to a moduli space;
\(\mathcal U\) transports metric, orientation, register, and boundary. Its
admissibility requires prolongation to every depth, commutation with
duality, descent of the spectrum to the quotient, and preservation of the domains and
codomains of the observables. These conditions convert an analogy of
ranks into a verifiable mathematical problem.

Strings appear after compactifying the eleven-dimensional sector;
branes appear as extended realizations coupled to the three-form and
its duals. The exceptional branch that reaches \(\Lambda_{24}\),
\(V^\natural\), and the Monster preserves the same enriched source through
another functor. The dimensional and exceptional interfaces are distinct
prolongations of a common genealogy, each with its own codomain.

\begin{resultado}
The perfection of \(G_{11}\), the polarization \(11\to10\), the isotropic
reduction \(26\to24\), the cubic screen of signature \((3,1)\) and the
radius involution are proven algebraic results. Together they form the
exact residence of ranks, quotients and dualities on which the interface
toward strings, branes and M-theory is built.
\end{resultado}

\begin{lecturafisica}
An extended history preserves position, winding, orientation, and
boundary within a single persistence. Compactification determines which
directions appear as macroscopic extension and which act as internal degrees.
The string reads that persistence on a worldsheet; the brane reads it on a
worldvolume; M-theory coordinates both readings through eleven-dimensional fields.
\end{lecturafisica}

\begin{frontera}
The realization is falsified if the map \(\mathfrak R_M\) alters the types of the screens, if the duality ceases to act unitarily on the spectrum, if the compactification loses prolongation, or if the fields fail to satisfy the action and supersymmetry. Physical closure requires tension, oscillators, amplitudes, controlled anomalies, and a coherent gravitational limit for each declared compactification.
\end{frontera}
